\chapter*{Abstract}

Diese Arbeit untersucht ein interaktives KI-Assistenzsystem im Vibe-Coding-Kontext im Rahmen einer standardisierten Coding-Challenge in einer kontrollierten Brownfield-Sandbox (BuyIn). Ziel ist es, Integrations- und Abgleichsarbeit im untersuchten Setting empirisch über Artefaktspuren und Interaktionsdaten beschreibbar zu machen, ohne Produktivitätsgewinne zu quantifizieren oder eine Rangordnung von Tools oder Modellen abzuleiten. Leitend sind drei Forschungsfragen zum Zusammenhang von Entwicklerkompetenz als Analyseperspektive mit beobachtbaren Umsetzungsständen, zur Rolle von Code- und Ausführungsartefakten für die Ergebnisbeschreibung sowie zu Chat-Interaktionsmustern im Arbeitsprozess.

Methodisch basiert die Studie auf einer standardisierten Coding-Challenge mit $N = 18$ Teilnehmenden, einer Timebox von 60 Minuten und einer standardisierten Entwicklungsumgebung. Die Auswertung nutzt eine artefaktgestützte Pipeline-Logik mit den Statuskategorien \textit{Attempted}, \textit{Implemented\textsubscript{static}} und \textit{Execution observed} und trianguliert diese mit Code- und Ausführungsartefakten, Chat-Exports und Survey-Merkmalen. Die Ergebnisdarstellung bleibt überwiegend deskriptiv; es werden keine inferenzstatistischen Gruppenvergleiche durchgeführt und keine kausalen Effekte beansprucht. Die Befunde sind auf das untersuchte Setting bezogen und sind nicht als allgemeingültige Integrations- oder Einführungsempfehlung zu verstehen.
